% !tp-example Basic usage
\documentclass{article}
\begin{document}
\wcTpTextToSpeech[lang=en-GB,show-text]{
  \wcScript[type=tp/txt]|
    Welcome to tp-components. This example uses the browser speech synthesis service.
  |
}
\end{document}

% !tp-example Javascript API
\documentclass{article}
\begin{document}
\wcScript[type=module]|
import "/components/text-to-speech/text-to-speech.js";

const example = document.createDocumentFragment();

const element1 = document.createElement("tp-text-to-speech");
element1.setAttribute("lang", "en-GB");
element1.setAttribute("show-text", "");
element1.append(document.createTextNode("\n  "));

const element2 = document.createElement("script");
element2.setAttribute("type", "tp/txt");
element2.textContent = "\n    Welcome to tp-components. This example uses the browser speech synthesis service.\n  ";
element1.append(element2);
element1.append(document.createTextNode("\n"));
example.append(element1);

document.body.append(example);
|
\end{document}

% !tp-example English voices
\documentclass{article}
\begin{document}
\wcTpTextToSpeech[id=voice-comparison,value=Welcome to tp-components. You can pause this reading and resume it whenever you want.,lang=en-US,show-text]{}

\wcScript[type=module,src=/docs/components/text-to-speech/examples/voice-comparison.js]||
\end{document}

% !tp-example French voices
\documentclass{article}
\begin{document}
\rawhtml|<tp-text-to-speech id="voice-comparison" value="Bienvenue dans tp-components. Vous pouvez mettre cette lecture en pause, puis la reprendre à votre rythme." lang="fr-FR" show-text=""></tp-text-to-speech>|

\wcScript[type=module,src=/docs/components/text-to-speech/examples/voice-comparison.js]||
\end{document}

% !tp-example Synchronized speech
\documentclass{article}
\begin{document}
Press Speak, then try Pause, Resume and Stop. Some voices do not report word boundaries; the full text remains visible in that case.


\wcTpTypewriting[id=spoken-text,word]{
  Welcome to \textbf{tp-components}. You can pause this reading and resume it whenever you want.

}

\wcTpTextToSpeech[for=spoken-text,lang=en]{}
\end{document}

% !tp-example Read an element
\documentclass{article}
\begin{document}
\wcArticle[id=reading-article]{
  Welcome to \textbf{tp-components}. This article stays visible while it is read aloud.

  You can pause the reading, resume it, or stop it whenever you want.

}

\wcTpTextToSpeech[for=reading-article,lang=en]{}
\end{document}

% !tp-example Pauses
\documentclass{article}
\begin{document}
\wcTpTextToSpeech[lang=en-GB,show-text,pause-between-lines=800]{
  \wcScript[type=tp/txt]|
    Look at the illustration. [pause=1500]
    A. A woman is placing [pause=300] a box on a table.
    B. Two people are boarding a train.
    C. A man is cleaning a window.
    D. Some chairs are being stacked.
  |
}
\end{document}

% !tp-example Dialogue
\documentclass{article}
\begin{document}
Voice names depend on your browser and system. This dialogue uses Daniel and Kate when available; otherwise it uses the component voice. Open Speech voice settings to see the available names for this language.


\wcTpTextToSpeech[lang=en-GB,show-text,pause-between-lines=400]{
  \wcScript[type=tp/txt]|[voice=Daniel] Hello! How are you?
[voice=Kate] Very well, thank you. [pause=600]
[voice=Daniel] Shall we begin?
[voice=default] End of the dialogue.|
}
\end{document}

% !tp-example Named triggers
\documentclass{article}
\begin{document}
\wcTpTextToSpeech[data-trigger-demo,lang=en-GB,show-text]{
  \wcScript[type=tp/txt]|Listen to the first idea. [trigger=clic] Now the second idea. [pause=500] [trigger=clic] Finished.|
}

\wcDiv[data-trigger-status,role=status]{No trigger yet.}

\wcScript[type=module,src=/docs/components/text-to-speech/examples/triggers.js]||
\end{document}

% !tp-example Automatic pauses
\documentclass{article}
\begin{document}
\wcTpTextToSpeech[lang=en-GB,show-text,pause-after-comma=200,pause-after-semicolon=400,pause-after-colon=400,pause-after-sentence=600,pause-per-blank-line=1000]{
  \wcScript[type=tp/txt]|First, listen to this idea; then consider the next one: punctuation sets the pace. Do you hear the pauses? Good!

A new paragraph follows one blank line.


Two blank lines add a longer break.|
}
\end{document}
